\documentclass[11pt]{article}
\usepackage{./UTIL/cw2}

\title{25WS72 Coursework 2: Architectural Design of Remote Control Unit for Remotely Operated Unmanned Vehicle}
\author{Matthew Morton - F537796}

\begin{document}

\begin{center}
    \textbf{25WS72: Systems Architecture Coursework 2}

    \textbf{Architectural Design of Remote Control Unit for Remotely Operated Unmanned Vehicle}

    \textbf{F537796}
\end{center}

\linespread{1.15}

\section*{Abstract}

Autonomy in marine applications is becoming increasingly popular, having developed significantly over the last few decades.
Whilst confidence in uncrewed vessels is growing, there are still questions around their safety and reliability.
This presents the need for a method of control for a human operator at certain stages of a mission.
One such phase is that where the vessel is within line of sight of the operator, such as in docking and undocking.
To meet this need, an initial architecture for a remote control unit (RCU) has been created.
This architecture aims to meet the core functional requirements set out by regulation such as the Maritime and Coastguard Agency's Workboat Code Edition 3 Annexe 2.
Such a remote controller is to be portable and usable in a harbour environment, and must provide control no worse than that of a coxswain aboard a vessel.
A structural, behavioural, and requirements model has been created in SysML 1.5, and traceability has been created between all three to communicate the architecture clearly to a variety of stakeholders.
This has resulted in a tablet-like device with physical human machine interface (HMI) controls that sits between the operator and the ROUV\@.
It converts human inputs to radiofrequency (RF) signals for control of the ROUV, and receives RF signals which it displays to the operator as information or alerts.
It utilises a hub and spoke structure centreing around a central processor, similar to most modern PCs.
This serves as a good base for further verification, validation, and design work to create a product ready for this space in the market.

\hrulefill

\tableofcontents

\hrulefill

\clearpage

\section{Introduction}

\subfile{SECTIONS/s1}

\clearpage

\section{Technical Results}

\subfile{SECTIONS/s2}

\clearpage

\section{Conclusion}

\subfile{SECTIONS/s3}

\clearpage

\printbibliography[heading=bibintoc, title=References]{}

\end{document}
